% Options for packages loaded elsewhere
\PassOptionsToPackage{unicode}{hyperref}
\PassOptionsToPackage{hyphens}{url}
\PassOptionsToPackage{dvipsnames,svgnames,x11names}{xcolor}
\PassOptionsToPackage{space}{xeCJK}
\documentclass[
]{article}
\usepackage{xcolor}
\usepackage[margin=2.5cm]{geometry}
\usepackage{amsmath,amssymb}
\setcounter{secnumdepth}{-\maxdimen} % remove section numbering
\usepackage{iftex}
\ifPDFTeX
  \usepackage[T1]{fontenc}
  \usepackage[utf8]{inputenc}
  \usepackage{textcomp} % provide euro and other symbols
\else % if luatex or xetex
  \usepackage{unicode-math} % this also loads fontspec
  \defaultfontfeatures{Scale=MatchLowercase}
  \defaultfontfeatures[\rmfamily]{Ligatures=TeX,Scale=1}
\fi
\usepackage{lmodern}
\ifPDFTeX\else
  % xetex/luatex font selection
  \setmainfont[]{DejaVu Sans}
  \setmonofont[]{DejaVu Sans Mono}
  \ifXeTeX
    \usepackage{xeCJK}
    \setCJKmainfont[]{Noto Sans CJK SC}
  \fi
  \ifLuaTeX
    \usepackage[]{luatexja-fontspec}
    \setmainjfont[]{Noto Sans CJK SC}
  \fi
\fi
% Use upquote if available, for straight quotes in verbatim environments
\IfFileExists{upquote.sty}{\usepackage{upquote}}{}
\IfFileExists{microtype.sty}{% use microtype if available
  \usepackage[]{microtype}
  \UseMicrotypeSet[protrusion]{basicmath} % disable protrusion for tt fonts
}{}
\makeatletter
\@ifundefined{KOMAClassName}{% if non-KOMA class
  \IfFileExists{parskip.sty}{%
    \usepackage{parskip}
  }{% else
    \setlength{\parindent}{0pt}
    \setlength{\parskip}{6pt plus 2pt minus 1pt}}
}{% if KOMA class
  \KOMAoptions{parskip=half}}
\makeatother
\usepackage{longtable,booktabs,array}
\usepackage{caption}
\captionsetup[table]{skip=6pt}
\newcounter{none} % for unnumbered tables
\usepackage{calc} % for calculating minipage widths
% Correct order of tables after \paragraph or \subparagraph
\usepackage{etoolbox}
\makeatletter
\patchcmd\longtable{\par}{\if@noskipsec\mbox{}\fi\par}{}{}
\makeatother
% Allow footnotes in longtable head/foot
\IfFileExists{footnotehyper.sty}{\usepackage{footnotehyper}}{\usepackage{footnote}}
\makesavenoteenv{longtable}
\setlength{\emergencystretch}{3em} % prevent overfull lines
\providecommand{\tightlist}{%
  \setlength{\itemsep}{0pt}\setlength{\parskip}{0pt}}
\usepackage{bookmark}
\IfFileExists{xurl.sty}{\usepackage{xurl}}{} % add URL line breaks if available
\urlstyle{same}
% fallback for those not using the hyperref driver hyperxmp:
\makeatletter
\@ifundefined{xmpquote}{\newcommand{\xmpquote}[1]{#1}}{}
\makeatother
\hypersetup{
  colorlinks=true,
  linkcolor={Maroon},
  filecolor={Maroon},
  citecolor={Blue},
  urlcolor={Blue},
  pdfcreator={LaTeX via pandoc}}

\author{}
\date{}

\begin{document}

\section{筑基篇课后习题 VIII：pat
原生群表示理论}\label{ux7b51ux57faux7bc7ux8bfeux540eux4e60ux9898-viiipat-ux539fux751fux7fa4ux8868ux793aux7406ux8bba}

\begin{quote}
\textbf{声明 (Declaration)}: 本文\textbf{不代表任何数学结论}
(non-mathematical-claim), 仅为基于 Lean 4
形式化的一种\textbf{实验观测报告} (experimental observation report)。

{\def\LTcaptype{none} % do not increment counter
\begin{longtable}[]{@{}
  >{\raggedright\arraybackslash}p{(\linewidth - 6\tabcolsep) * \real{0.2500}}
  >{\raggedright\arraybackslash}p{(\linewidth - 6\tabcolsep) * \real{0.2500}}
  >{\raggedright\arraybackslash}p{(\linewidth - 6\tabcolsep) * \real{0.2500}}
  >{\raggedright\arraybackslash}p{(\linewidth - 6\tabcolsep) * \real{0.2500}}@{}}
\toprule\noalign{}
\begin{minipage}[b]{\linewidth}\raggedright
习题
\end{minipage} & \begin{minipage}[b]{\linewidth}\raggedright
主题
\end{minipage} & \begin{minipage}[b]{\linewidth}\raggedright
Lean 形式化
\end{minipage} & \begin{minipage}[b]{\linewidth}\raggedright
DOI
\end{minipage} \\
\midrule\noalign{}
\endhead
\bottomrule\noalign{}
\endlastfoot
exercise\_i\_8 & pat 原生群表示理论（R165） & PatRepresentation.lean &
10.5281/zenodo.21916858 \\
\end{longtable}
}
\end{quote}

\textbf{Exercise VIII for the Foundation-Building Chapter: PAT-Native
Group Representation Theory}

\begin{quote}
习题定位：筑基篇课后习题系列第八题。用户指令：开始尝试对群表示理论进行
Lean 形式化，但\textbf{不抄 mathlib}------pat
原生定义。全部新定理只锚筑基篇定理（R138/R141/R143/R161），不用外部引理，不
import mathlib RepresentationTheory。
\end{quote}

\emph{2026-08-13 · Internal research exercise · Lean 4 / mathlib v4.32.2
· 7 new theorems PROVED no sorry · 算器神魂论筑基篇课后习题 VIII }

\begin{center}\rule{0.5\linewidth}{0.5pt}\end{center}

\subsection{习题陈述}\label{ux4e60ux9898ux9648ux8ff0}

\begin{quote}
对群表示理论进行 Lean 形式化，pat 原生（用户纠正：不抄
mathlib）。唯一论点：\textbf{群的 pat 原生表示 =
相位表示------每个群元素映射到单位圆上的相位
exp(i·θ(g))，保持群结构；每个元素与其逆元构成互锁对（ρ(g)·ρ(g⁻¹) =
1）。}
\end{quote}

\begin{center}\rule{0.5\linewidth}{0.5pt}\end{center}

\subsection{零、总论：为什么群表示在 pat
里是相位表示}\label{ux96f6ux603bux8bbaux4e3aux4ec0ux4e48ux7fa4ux8868ux793aux5728-pat-ux91ccux662fux76f8ux4f4dux8868ux793a}

mathlib 的 RepresentationTheory 定义表示为
\texttt{G\ →*\ V\ →ₗ{[}k{]}\ V}（群到线性算子的同态）------这是''抄现成的''。pat
原生视角完全不同：

\textbf{pat 框架的表示 = 相位表示}（RulerPhase：相位差 =
方向；R138：相位关系锁定）： - 群元素 g ↦ 单位圆相位 exp(i·θ(g)) -
同态性：θ(gh) = θ(g) + θ(h)（相位差可加） - 单位元：ρ(1) = exp(0) = 1 -
\textbf{互锁对}：ρ(g⁻¹) = ρ(g)⁻¹（逆元 = 反向相位），ρ(g)·ρ(g⁻¹) = 1

这不是''抄一个更弱的结构''------这是\textbf{从互锁原则推出的表示理论}：表示的结构单元不是线性算子，而是互锁对
\{ρ(g), ρ(g)⁻¹\}（R136 ②③：方向必须成对声明）。

\begin{center}\rule{0.5\linewidth}{0.5pt}\end{center}

\subsection{一、相位表示：群 →
单位圆（RulerPhase/R141）}\label{ux4e00ux76f8ux4f4dux8868ux793aux7fa4-ux5355ux4f4dux5706rulerphaser141}

\textbf{★核心：群 G 的相位表示 = 相位函数 θ : G → ℝ 保持群结构；表示值
ρ(g) = exp(i·θ(g)) 在单位圆上。}

\subsubsection{论证}\label{ux8bbaux8bc1}

\begin{enumerate}
\def\labelenumi{\arabic{enumi}.}
\tightlist
\item
  \textbf{相位差 =
  方向}（RulerPhase）：群的每个元素映射到一个相位（方向）。
\item
  \textbf{表示值在单位圆}（R141）：‖exp(i·θ(g))‖ =
  1------单位根相位圆；U(1) 规范结构的单相位圆。
\item
  \textbf{单位元 = 相位 0}（exp\_zero/R143）：ρ(1) = exp(0) =
  1------对称对还原到 1。
\end{enumerate}

\subsubsection{形式化（PatRepresentation.lean，3
定理）}\label{ux5f62ux5f0fux5316patrepresentation.lean3-ux5b9aux7406}

\begin{itemize}
\tightlist
\item
  \texttt{phaseRep\_value}（def）：ρ(g) = exp(i·θ(g))。
\item
  \texttt{phaseRep\_on\_circle}：‖exp(i·θ(g))‖ = 1------表示值在单位圆。
\item
  \texttt{phaseRep\_one}：exp(i·0) = 1------单位元 = 相位 0。
\end{itemize}

\textbf{验证}：0 sorry。锚定 Complex.norm\_exp\_ofReal\_mul\_I / simp。

\begin{center}\rule{0.5\linewidth}{0.5pt}\end{center}

\subsection{二、同态性与互锁对（R138/R143/R136
②③）}\label{ux4e8cux540cux6001ux6027ux4e0eux4e92ux9501ux5bf9r138r143r136-ux2461ux2462}

\textbf{★核心：相位差可加 ⟹ 同态；逆元 = 反向相位 ⟹ 互锁对还原
ρ(g)·ρ(g⁻¹) = 1。}

\subsubsection{论证}\label{ux8bbaux8bc1-1}

\begin{enumerate}
\def\labelenumi{\arabic{enumi}.}
\tightlist
\item
  \textbf{同态性}（R138：相位锁定后相位差可加）：θ(gh) = θ(g) + θ(h) ⟹
  ρ(gh) = ρ(g)·ρ(h)（exp\_add）------表示的群同态性质。
\item
  \textbf{逆元 = 反向相位}（R136 ②③：方向必须成对声明）：θ(g⁻¹) = -θ(g)
  ⟹ ρ(g⁻¹) = ρ(g)⁻¹（exp\_neg）------互锁对的逆元侧。
\item
  \textbf{★互锁对还原}（R143：对称对还原到 1）：ρ(g)·ρ(g⁻¹) =
  exp(iθ)·exp(-iθ) = 1------每个群元素与其逆元构成互锁对，组合还原到
  1。\textbf{这是''互锁 = 成对''在群表示中的对应。}
\end{enumerate}

\subsubsection{形式化（PatRepresentation.lean，3
定理）}\label{ux5f62ux5f0fux5316patrepresentation.lean3-ux5b9aux7406-1}

\begin{itemize}
\tightlist
\item
  \texttt{phaseRep\_mul}：θ(gh) = θ(g)+θ(h) ⟹ ρ(gh) =
  ρ(g)·ρ(h)------同态性。
\item
  \texttt{phaseRep\_inv}：θ(g⁻¹) = -θ(g) ⟹ ρ(g⁻¹) = ρ(g)⁻¹------逆元 =
  反向相位。
\item
  \texttt{phaseRep\_pair\_reduces}：ρ(g)·ρ(g⁻¹) = 1------★互锁对还原。
\end{itemize}

\textbf{验证}：0 sorry。锚定 exp\_add / exp\_neg / exp\_zero + ring。

\begin{center}\rule{0.5\linewidth}{0.5pt}\end{center}

\subsection{三、★规范群 = k
个独立相位参数（R161）}\label{ux4e09ux89c4ux8303ux7fa4-k-ux4e2aux72ecux7acbux76f8ux4f4dux53c2ux6570r161}

\textbf{★核心：U(1) = 1 相位参数（1 对互锁 = 2 互锁），SU(2) = 2
相位参数（2 对 = 4 互锁 = S³），SU(3) = 4 相位参数（4 对 = 8
互锁）------k 个独立相位参数各自给互锁对。}

\subsubsection{论证}\label{ux8bbaux8bc1-2}

\begin{enumerate}
\def\labelenumi{\arabic{enumi}.}
\tightlist
\item
  \textbf{k 对独立互锁自洽}（R161
  k\_pairs\_independent\_interlock）：任意 k 个独立相位参数 θ : Fin k →
  ℝ，每对 exp(iθⱼ)·exp(-iθⱼ) = 1。
\item
  \textbf{规范群对应}：U(1) = 1 参数（单相位圆 R138）；SU(2) = 2
  参数（S³，R149/R154，习题 VII 量子观测）；SU(3) = 4 参数 = 8 互锁。
\item
  \textbf{诚实边界}：生成元与互锁对的精确对应需群表示理论（CONJECTURE
  层）------本定理锚定的是 k 对独立互锁自洽（R161），非群同构证明。
\end{enumerate}

\subsubsection{形式化（PatRepresentation.lean，1
定理）}\label{ux5f62ux5f0fux5316patrepresentation.lean1-ux5b9aux7406}

\begin{itemize}
\tightlist
\item
  \texttt{gauge\_k\_phase\_params}：任意 k
  个独立相位参数，每对互锁自洽------规范群结构。
\end{itemize}

\textbf{验证}：0 sorry。锚定 k\_pairs\_independent\_interlock。

\begin{center}\rule{0.5\linewidth}{0.5pt}\end{center}

\subsection{四、全景（组合定理）}\label{ux56dbux5168ux666fux7ec4ux5408ux5b9aux7406}

\textbf{★核心：表示值在单位圆 ∧ 单位元 = 相位 0 ∧ 同态性 ∧
互锁对还原------pat 原生群表示理论的核心结构。}

\subsubsection{形式化（PatRepresentation.lean，1
定理）}\label{ux5f62ux5f0fux5316patrepresentation.lean1-ux5b9aux7406-1}

\begin{itemize}
\tightlist
\item
  \texttt{phaseRep\_perspective}：‖ρ(g)‖ = 1 ∧ ρ(gh) = ρ(g)·ρ(h) ∧
  ρ(g)·ρ(g⁻¹) = 1------全景。
\end{itemize}

\textbf{验证}：0 sorry。合取项分别锚 phaseRep\_on\_circle /
phaseRep\_mul / phaseRep\_pair\_reduces。

\begin{center}\rule{0.5\linewidth}{0.5pt}\end{center}

\subsection{五、习题解答总结}\label{ux4e94ux4e60ux9898ux89e3ux7b54ux603bux7ed3}

{\def\LTcaptype{none} % do not increment counter
\begin{longtable}[]{@{}
  >{\raggedright\arraybackslash}p{(\linewidth - 6\tabcolsep) * \real{0.2500}}
  >{\raggedright\arraybackslash}p{(\linewidth - 6\tabcolsep) * \real{0.2500}}
  >{\raggedright\arraybackslash}p{(\linewidth - 6\tabcolsep) * \real{0.2500}}
  >{\raggedright\arraybackslash}p{(\linewidth - 6\tabcolsep) * \real{0.2500}}@{}}
\toprule\noalign{}
\begin{minipage}[b]{\linewidth}\raggedright
论点
\end{minipage} & \begin{minipage}[b]{\linewidth}\raggedright
内容
\end{minipage} & \begin{minipage}[b]{\linewidth}\raggedright
锚到
\end{minipage} & \begin{minipage}[b]{\linewidth}\raggedright
定理
\end{minipage} \\
\midrule\noalign{}
\endhead
\bottomrule\noalign{}
\endlastfoot
相位表示 & 群 → 单位圆 exp(i·θ(g))；值在单位圆；单位元 = 1 &
RulerPhase/R141/R143 & phaseRep\_on\_circle / phaseRep\_one \\
同态性 & θ(gh) = θ(g)+θ(h) ⟹ ρ(gh) = ρ(g)·ρ(h) & R138 & phaseRep\_mul \\
★互锁对 & ρ(g)·ρ(g⁻¹) = 1（逆元 = 反向相位） & R143/R136②③ &
phaseRep\_inv / phaseRep\_pair\_reduces \\
规范群 & k 独立相位参数 = k 对互锁 = 2k 互锁 & R161 &
gauge\_k\_phase\_params \\
★全景 & 单位圆 + 同态 + 互锁对 & 上述全部 & phaseRep\_perspective \\
\end{longtable}
}

\textbf{习题的回答（一句话）}：\textbf{群的 pat 原生表示 =
相位表示------群元素映射到单位圆相位
exp(i·θ(g))，保持群结构；每个元素与其逆元构成互锁对（ρ(g)·ρ(g⁻¹) =
1，R143 对称对还原）------表示的结构单元是互锁对，不是线性算子（不抄
mathlib，pat 原生）。}

\textbf{对 R164 强弱力观测的支撑}：相位表示给出 U(1) = 1
相位参数（单相位圆）、SU(2) = 2 参数（4 互锁 = S³）、SU(3) = 4 参数（8
互锁）的结构------R164 的''规范群 = 互锁对数''从 CONJECTURE 推进到''k
对独立互锁自洽已 PROVED + 精确生成元对应仍 CONJECTURE''。

\begin{center}\rule{0.5\linewidth}{0.5pt}\end{center}

\subsection{六、相关对照}\label{ux516dux76f8ux5173ux5bf9ux7167}

{\def\LTcaptype{none} % do not increment counter
\begin{longtable}[]{@{}
  >{\raggedright\arraybackslash}p{(\linewidth - 4\tabcolsep) * \real{0.3333}}
  >{\raggedright\arraybackslash}p{(\linewidth - 4\tabcolsep) * \real{0.3333}}
  >{\raggedright\arraybackslash}p{(\linewidth - 4\tabcolsep) * \real{0.3333}}@{}}
\toprule\noalign{}
\begin{minipage}[b]{\linewidth}\raggedright
文献
\end{minipage} & \begin{minipage}[b]{\linewidth}\raggedright
对应内容
\end{minipage} & \begin{minipage}[b]{\linewidth}\raggedright
备注
\end{minipage} \\
\midrule\noalign{}
\endhead
\bottomrule\noalign{}
\endlastfoot
\textbf{Weyl, H.} 1939. \emph{The Classical Groups} & 群表示理论经典 &
pat 原生对照（不抄） \\
\textbf{R138/R141/R143/R161}（筑基篇） & 相位锁定/单位根圆/对称对还原/k
对互锁 & 本框架定理，非外部引理 \\
\end{longtable}
}

\begin{quote}
诚实边界：本习题自建 pat 原生表示（相位表示），不使用 mathlib
RepresentationTheory（用户纠正''不抄''）；与经典群表示的关系为结构对照（非替代）。
\end{quote}

\begin{center}\rule{0.5\linewidth}{0.5pt}\end{center}

\emph{筑基篇课后习题 VIII · 2026-08-13 · Lean 4 / mathlib v4.32.2 · 7
theorems PROVED no sorry（PatRepresentation.lean）· 配套 Lean
形式化：formal/Formal/Toolkit/PatRepresentation.lean（快照见
../lean/PatRepresentation.lean）}

\end{document}
